\section{Limitations and conclusion}

Several limitations restrict inference. GEDI footprints of approximately 25 m are represented on a 10 m grid, introducing a support mismatch that is especially relevant near boundaries. The 2020--2021 composites may differ temporally from individual lidar observations, and the mangrove mask can introduce inclusion or omission errors. Nonoverlapping extraction leaves uncovered margins, while sparse-reference filtering limits the range of patches used for assessment. Evaluation uses GEDI alone, with only 163 retained test cells across five patches; independent field or airborne lidar validation is unavailable in the supplied experiments. Run-summary inconsistencies must also be resolved before definitive numerical reporting.

We investigated optical--SAR fusion for mangrove canopy-height regression in Cà Mau using sparse GEDI RH98 supervision and spatially separated evaluation. The supplied aggregate results favor dual-encoder fusion on test over optical-only and early-fusion models, while validation gains are small and metric dependent. The study supports a cautious application-level conclusion: modality-specific encoding is a promising option for this regional dataset, but its advantage, dense-map accuracy, and transferability require further verification. Immediate priorities are reconciliation of run logs, completion of height-dependent diagnostics, and independent validation at an appropriate spatial support.
